\scanpage{184}

\uline{Remarques} a) Si $X$ est compact, on voit qu'il n'y a qu'un nb fini de \deleted{feuilles} \inserted{strates} (qui sont loc. compactes) --- on s'intéresse surtout aux \deleted{feuilletages} \inserted{stratifications} avec un nb fini de strates --- mais bien sûr cette hyp. n'est pas stable par passage à un ouvert de $X$.

b) La formation de $D_{\Delta_0}^{*}$ au dessus de $X$ commute à la restriction à un ouvert de $X$.

c) Si la \deleted{\unsure{partition}} \inserted{structure} de \unsure{strat.} est définie par un diviseur à croisements normaux sur un espace analytique, alors on a un homéo.
\[
(298)\qquad D_{\Delta_0}^{*}\simeq\coprod_{d\in\mathbb N}D_d^{*}
\]
donc la notation $D_{\Delta_0}^{*}$ introduite ici coïncide avec celle des \deleted{\unsure{la}} deux sections précédentes.

\margin{183}
